\documentclass[11pt]{article}
\usepackage[margin=1in]{geometry}
\usepackage{amsmath,amssymb,hyperref}
\title{Fast Archipelago v1.3\\One-shot fast agreement with certified slow fallback}
\author{Protocol revision and implementation specification}
\date{4 October 2026}
\setlength{\emergencystretch}{2em}
\begin{document}
\maketitle
\begin{abstract}
This revision replaces the repeated per-rank fast/recovery construction in v1.2.
Each checkpoint has one immutable fast-vote attempt. Recovery certificates
supply externally valid initial inputs to a fresh ordinary Archipelago instance.
The slow instance runs its own R, A and B phases; recovery never overrides its
R result. Fast evidence remains actionable after slow execution starts.
The fast/recovery composition argument is given below. End-to-end termination
remains conditional on the stated ordinary slow-engine and dissemination
obligations; this document does not claim those implementation obligations are
already discharged.
\end{abstract}

\section{Model and domains}
Fix a checkpoint instance $I$ binding protocol version, session, epoch,
predecessor checkpoint and the exact equal-weight committee and fault budgets.
There are $n=3f+2p+1$ identities, at most $f$ Byzantine identities, and
unforgeable signatures. Correct identities never change their first vote in $I$.
Let
\[
 F=n-p,\qquad Q=n-f-p=2f+p+1,\qquad P=n-f=2f+2p+1.
\]
These are fast/recovery thresholds, not the slow engine's quorum thresholds.
Values have a deterministic total order and must satisfy the application validity
predicate, including reconstructed payload availability and predecessor binding.
A hash alone does not establish application validity.

For the concrete fallback in this revision, choose the lexicographically smallest
$3f+1$ public keys from the fast committee, fixed for $I$. Define $I_s$ by hashing
an independent slow-instance domain, $I$, $f$ and the ordered selected keys.
The slow committee has at most $f$ Byzantine identities, quorum $q_s=2f+1$
and witness threshold $w_s=f+1$. Fast members outside this committee still
serve fast votes, recovery proofs and application data. Changing this selection
rule requires a protocol version change.

Every signed object binds a message-kind domain, its instance, signer, value
and relevant phase/rank/request identifiers. Fast attempt number is always zero;
slow rank zero is a distinct namespace. v1.2 signatures are not v1.3 signatures.
A slow proposal signature binds its value and slow instance; its attached RC
is independently verified. Different valid RC witnesses for the same value are
interchangeable. Subsequent slow phase signatures must bind their phase-specific
proof references under the ordinary engine's rules.


\section{Finite application admission and initial R evidence}
Before releasing FIRST, the RAI binding authenticates a candidate object containing
$I$, its proposer and the complete epoch value. Validation reconstructs exactly
$P=n-f$ distinct committee reports, checks both roots and their signed committee
and predecessor context, and recomputes BuildState. A signature alone is not
application validation.

A candidate becomes an admissible initial value only with $Q$ distinct signed
endorsements of that exact candidate object and proposer. Each correct endorser
locks one candidate per $(I,\mathrm{proposer})$ for the process lifetime. Two
admission quorums for one proposer intersect in at least $2Q-n=f+1$ identities;
at least one is correct and cannot endorse both. Hence at most one admitted
candidate exists per proposer, and there are at most $n$ initial values. Correct
proposers must release one validated candidate and dissemination must make its
payload available. Admission is an additional exchange before the optimistic
FIRST attempt; latency measurements must include it when measuring the complete
checkpoint operation. Disk/restart safety for endorsement locks is deferred.

For empty-candidate recovery the implementation adds an independent initial
max-register exchange, in a domain separate from FIRST and the RC-backed slow
instance. An initial R request names an admitted candidate. A correct responder
incorporates the request into its retained maximum and signs a response binding
$I$, the request digest and its maximal admitted candidate. An initial R
certificate contains $Q$ distinct such responses to the same authenticated
request. Verification checks signatures, admitted origins, and that each response
is at least the triggering input, and recomputes their maximum. The cache stores
only verified certificates; the FIRST snapshot's maximum is never substituted.
Candidate-admission quorums witness origin availability in this initial exchange;
these are not the later slow engine's R-response witnesses. This specifies the
initial max-register refinement used by the application adapter. It does not
supply the ordinary slow synchronizer or discharge its termination assumptions.

\section{One initial fast attempt}
The initial fast R exchange retains the eligible-request/response and witness
checks of the fast/recovery front end. A correct replica authenticates a valid
initial proposal and releases at most one signed $\mathrm{FIRST}(I,0,v)$.
Later information can change an ordinary R response, but cannot change FIRST.
The immutable record must be installed atomically before releasing the signature.
There are no FIRST votes at slow ranks, including later slow ranks.

A fast certificate $FC_0(v)$ contains exactly $F$ distinct valid FIRST signatures
for $v$. Receipt of such a certificate permits deciding $v$ at any time,
including while the slow engine is running. Correct replicas relay certificates
and dependencies; starting the fallback does not disable this listener.

The implementation currently retains locks only in memory, at the user's
request. Therefore its safety scope excludes restart with lost locks followed
by signing again in the same instance. Durable restart support is deferred.

\section{Recovery and admission}
Let $S$ contain one validated FIRST record per identity, with $Q\leq m=|S|\leq n$.
Define
\[
 \mathrm{Cand}(S)=\{v:\mathrm{count}_S(v)\geq m-f-p\}.
\]
An equivocating identity contributes at most one record to a snapshot. A receiver
may nevertheless retain alternate signed records for detecting a fast certificate.

A resolved certificate $RC_0(x)$ contains $I$, attempt zero, $S$, $x$, and one of:
\begin{enumerate}
\item Singleton evidence: $\mathrm{Cand}(S)=\{x\}$.
\item Empty evidence: $\mathrm{Cand}(S)=\varnothing$ and independently verified
initial ordinary R evidence whose recomputed maximum is $x$.
\end{enumerate}
An ambiguous candidate set does not resolve: neither taking its maximum nor
replacing old first votes is allowed. Missing dependencies are fetched and retried;
invalid dependencies do not count. Empty-case R evidence must include all required
ancestry and eligibility witnesses, not merely a caller-supplied maximum.

Any member of the slow committee may start the fresh slow instance with an
initial proposal $x$ accompanied by $RC_0(x)$. All initial slow proposals,
including Byzantine proposals, must satisfy this external validity predicate.
Validation is independent of the verifier's local recovery snapshot. A valid
certificate does not rewrite that verifier's first vote.

\section{Timing and ordinary slow ranks}
There is a causal boundary: recovery evidence must exist before its value enters
slow consensus. There is no wait for a timeout, failed-fast proof, all fast votes,
or abandonment of the fast attempt. Slow work and delivery of late initial fast
evidence therefore overlap.

The slow instance starts with fresh responder, witness and synchronizer state.
Its ordinary rank structure is:
\begin{enumerate}
\item R: exchange legal requests, their witnesses and eligible responses;
select the ordinary R result $r$ from a valid slow quorum and retained local R state.
When the retained maximum exceeds the quorum maximum, its signed record and
origin witnesses must accompany the selection proof.
\item A: propose $r$ with its slow R evidence. Verify ordinary A request,
response and complementing-witness rules; produce the ordinary A result.
\item B: propose that A result with its proof; verify B evidence and either
commit or adopt according to the ordinary B evaluator.
\item On adoption, advance via the ordinary synchronizer and legal preceding-B
ancestry. Do not introduce a new fast attempt or replace the next R result by RC.
\end{enumerate}
Request collection barriers, response updates, W admission and Q eligibility
are internal rules of the chosen ordinary engine, not interchangeable messages
or additional consensus decisions. Their exact source-faithful implementation
and recursive proof validation remain required work.

In particular, if the initial external input is $v$ but slow R validly selects
$w$, A must propose $w$. The initial $RC_0(v)$ is not permission to override R.
Every value introduced later must have ordinary legal ancestry back to an
externally valid initial input. Neither an old fast certificate nor an arbitrary
hash substitutes for preceding slow B evidence.

\section{Safety of the composition}
\paragraph{Fast uniqueness.}
Two fast signer sets intersect in at least $2F-n=n-2p=3f+1>f$ identities.
Some intersecting identity is correct and cannot first-vote for two values.
Thus conflicting fast certificates cannot exist.

\paragraph{Hidden-fast preservation.}
Suppose $FC_0(v)$ exists, even if initially hidden. At least $F-f$ correct
identities first-voted $v$. A snapshot of size $m$ includes at least
\[
 (F-f)+m-n=m-f-p
\]
of them. Hence $v\in\mathrm{Cand}(S)$ for every valid recovery snapshot.
An empty certificate is impossible, and every resolved singleton certificate
selects $v$. This also applies to a fast certificate assembled later: a correct
identity that first-voted another value cannot subsequently join it for $v$.

\paragraph{Fast/slow agreement.}
Assume the slow engine enforces external validity on every initial input,
preserves legal ancestry, satisfies agreement, and decides $v$ when every
externally valid input equals $v$. If $FC_0(v)$ exists, hidden-fast preservation
restricts every admissible slow input to $v$, so every slow decision is $v$.
Together with fast uniqueness and slow agreement, decisions agree across paths.
These are explicit requirements on the ordinary engine, not properties established
by the current pure A/B helper functions alone.

\section{Recovery progress and conditional termination}
At $m\geq P$, two candidates would require at least $2(m-f-p)>m$ votes, since
$m\geq2f+2p+1$. Therefore there is at most one candidate. This establishes
resolvability of a sufficiently large snapshot; in the empty case it still
requires available, independently valid initial R evidence.

Assume reliable eventual dissemination of correct initial FIRST records and
required payloads/proofs, eventual production of a valid initial R certificate
when needed, and an ordinary slow engine satisfying termination under the
execution's scheduling and synchronizer assumptions. Once $P$ correct identities
participate, recovery resolves; correct slow members receive admissible proposals
and the slow engine terminates, or a fast certificate already decides the instance.
Thus the composition terminates under these premises.

For static temporary suspensions, the selected cohort may be unavailable for a
finite interval and later resume; this is not permanent absence. After the last
return, the fixed slow committee again has the ordinary at-most-$f$ Byzantine
fault model. A proof must show that the slow synchronizer recovers from the
arbitrary prefix and reaches its progress condition. Static identities alone do
not imply such a suffix if suspension intervals can recur forever. This revision
does not claim termination during arbitrary recurring or rotating suspensions,
nor does it assume all responses become identical merely because values are ordered.

\section{Why the repeated construction was withdrawn}
In v1.2, recovery could repeatedly replace an ordinary R result $w$ by $v$ before
A, while an adversary retained a privately assembled preceding adoption proof
for $w$. Mutable responder evidence permits earlier certificates to remain valid:
a later stronger certificate does not revoke signatures already released.
Consequently adoption of $v$ and adoption of $w$ can coexist without any commit.
The recurring trace in the accompanying investigation reproduces this gap even
with a clean suffix. A request barrier or a single update per rank is not, by
itself, a proof that a correct first-vote quorum eventually forms.

The repair removes the repeated override: RC is initial external validity for a
fresh slow instance, and ordinary R controls ordinary A at every slow rank.
This removes the exhibited recovery-induced cycle. It is not a substitute for
proving the slow engine's own agreement, validity, synchronizer and termination.
The former unrestricted adoption-persistence and unconditional composed-liveness
claims are not retained.

\section{Implementation contract and remaining gates}
Phase 1 now enforces attempt-zero fast signing and certificate verification,
versioned fast domains, deterministic slow membership/instance binding and signed
RC-backed initial slow proposals. Tests exercise rejection before signing,
recovery verification, replay isolation and late fast completion after fallback
admission. Existing A/B and witness helpers remain pure components.

The component implementation now includes signed slow requests, bounded recursive
ancestry checks and an in-memory proposer loop, including retained R selection.
Request proofs require weak complementing witnesses; quorum eligibility requires
strong witnesses. Component tests cover adoption and subsequent commitment.
Mutable responder services and digest-based proof fetch/relay now run in a
transport-independent replica. Four-replica simulations cover split inputs,
missing certificates and service to a returning replica after decision.
Typed-message admission and historical-response replay now have configurable byte
budgets; replay also has a count limit and rotates fairly without evicting evidence.
These are local accounting controls, not wire framing, total-storage limits or
round synchronization.
Finite initial candidates now have quorum admission and an independent initial R
certificate verifier. Transport codecs, a committee-wide driver and local
application validation are component implementations. Still required: admission
and application-evidence dissemination, retry scheduling and a round synchronizer,
live checkpoint service integration, adversarial timing coverage and Gate B
end-to-end checks. Aggregate retained-storage limits and disk persistence are
deferred; the current nanospam baseline keeps all nodes online. The one-shot interface tests are not an end-to-end termination proof.
Disk persistence remains deferred. No production checkpoint service is installed
by this revision.

The editable source is \texttt{doc/fast-archipelago/fast-archipelago.tex};
\texttt{rebuild.py} compiles the complete document. Historical v1.2 analyses and
replacement-page source remain separate investigation artifacts, not normative
parts of v1.3.
\end{document}
